\documentclass[10pt,a4paper]{amsart}
\usepackage{amsmath,amssymb,amsthm,graphicx}
% For commutative diagrams we use Xy-pic
\usepackage[all]{xy}

% We use 2cell for 2-commutative diagrams.
\xyoption{2cell}
% Theorem environments.
%
\theoremstyle{plain}
\newtheorem{theorem}[subsection]{Theorem}
\newtheorem{proposition}[subsection]{Proposition}
\newtheorem{lemma}[subsection]{Lemma}

\theoremstyle{definition}
\newtheorem{definition}[subsection]{Definition}
\newtheorem{example}[subsection]{Example}
\newtheorem{exercise}[subsection]{Exercise}
\newtheorem{situation}[subsection]{Situation}

\theoremstyle{remark}
\newtheorem{remark}[subsection]{Remark}
\newtheorem{remarks}[subsection]{Remarks}

\numberwithin{equation}{subsection}

% Macros
%
\def\lim{\mathop{\mathrm{lim}}\nolimits}
\def\colim{\mathop{\mathrm{colim}}\nolimits}
\def\Spec{\mathop{\mathrm{Spec}}}
\def\Hom{\mathop{\mathrm{Hom}}\nolimits}
\def\Ext{\mathop{\mathrm{Ext}}\nolimits}
\def\SheafHom{\mathop{\mathcal{H}\!\mathit{om}}\nolimits}
\def\SheafExt{\mathop{\mathcal{E}\!\mathit{xt}}\nolimits}
\def\Sch{\mathit{Sch}}
\def\Mor{\mathop{\mathrm{Mor}}\nolimits}
\def\Ob{\mathop{\mathrm{Ob}}\nolimits}
\def\Sh{\mathop{\mathit{Sh}}\nolimits}
\def\NL{\mathop{N\!L}\nolimits}
\def\CH{\mathop{\mathrm{CH}}\nolimits}
\def\proetale{{pro\text{-}\acute{e}tale}}
\def\etale{{\acute{e}tale}}
\def\QCoh{\mathit{QCoh}}
\def\Ker{\mathop{\mathrm{Ker}}}
\def\Im{\mathop{\mathrm{Im}}}
\def\Coker{\mathop{\mathrm{Coker}}}
\def\Coim{\mathop{\mathrm{Coim}}}

% Boxtimes
%
\DeclareMathSymbol{\boxtimes}{\mathbin}{AMSa}{"02}

%
% Macros for moduli stacks/spaces
%
\def\QCohstack{\mathcal{QC}\!\mathit{oh}}
\def\Cohstack{\mathcal{C}\!\mathit{oh}}
\def\Spacesstack{\mathcal{S}\!\mathit{paces}}
\def\Quotfunctor{\mathrm{Quot}}
\def\Hilbfunctor{\mathrm{Hilb}}
\def\Curvesstack{\mathcal{C}\!\mathit{urves}}
\def\Polarizedstack{\mathcal{P}\!\mathit{olarized}}
\def\Complexesstack{\mathcal{C}\!\mathit{omplexes}}
% \Pic is the operator that assigns to X its picard group, usage \Pic(X)
% \Picardstack_{X/B} denotes the Picard stack of X over B
% \Picardfunctor_{X/B} denotes the Picard functor of X over B
\def\Pic{\mathop{\mathrm{Pic}}\nolimits}
\def\Picardstack{\mathcal{P}\!\mathit{ic}}
\def\Picardfunctor{\mathrm{Pic}}
\def\Deformationcategory{\mathcal{D}\!\mathit{ef}}
\usepackage{cleveref}
\crefname{originallemma}{original lemma}{original lemmas}
\Crefname{originallemma}{Original lemma}{Original lemmas}
\makeatletter
\newlabel{algebra-lemma-pushdown-module}{{02M0}{source}{}{sourcepointer-1}{}}
\newlabel{algebra-lemma-pushdown-module@cref}{{[originallemma][0][]02M0}{[1][0][]source}{}{}{}}
\makeatother
\crefname{originallemma}{original lemma}{original lemmas}
\Crefname{originallemma}{Original lemma}{Original lemmas}
\makeatletter
\newlabel{coherent-lemma-coherent-support-closed}{{01Y5}{source}{}{sourcepointer-2}{}}
\newlabel{coherent-lemma-coherent-support-closed@cref}{{[originallemma][0][]01Y5}{[1][0][]source}{}{}{}}
\makeatother
\crefname{originallemma}{original lemma}{original lemmas}
\Crefname{originallemma}{Original lemma}{Original lemmas}
\makeatletter
\newlabel{coherent-proposition-proper-pushforward-coherent}{{02O5}{source}{}{sourcepointer-3}{}}
\newlabel{coherent-proposition-proper-pushforward-coherent@cref}{{[originallemma][0][]02O5}{[1][0][]source}{}{}{}}
\makeatother
\crefname{originallemma}{original lemma}{original lemmas}
\Crefname{originallemma}{Original lemma}{Original lemmas}
\makeatletter
\newlabel{lemma-cycle-closed-coherent}{{02QY}{source}{}{sourcepointer-4}{}}
\newlabel{lemma-cycle-closed-coherent@cref}{{[originallemma][0][]02QY}{[1][0][]source}{}{}{}}
\makeatother
\crefname{originallemma}{original lemma}{original lemmas}
\Crefname{originallemma}{Original lemma}{Original lemmas}
\makeatletter
\newlabel{lemma-multiplicity-finite}{{02QT}{source}{}{sourcepointer-5}{}}
\newlabel{lemma-multiplicity-finite@cref}{{[originallemma][0][]02QT}{[1][0][]source}{}{}{}}
\makeatother
\crefname{originallemma}{original lemma}{original lemmas}
\Crefname{originallemma}{Original lemma}{Original lemmas}
\makeatletter
\newlabel{morphisms-lemma-generically-finite}{{02NW}{source}{}{sourcepointer-6}{}}
\newlabel{morphisms-lemma-generically-finite@cref}{{[originallemma][0][]02NW}{[1][0][]source}{}{}{}}
\makeatother
\crefname{originallemma}{original lemma}{original lemmas}
\Crefname{originallemma}{Original lemma}{Original lemmas}
\makeatletter
\newlabel{morphisms-lemma-quasi-compact-scheme-theoretic-image}{{01R8}{source}{}{sourcepointer-7}{}}
\newlabel{morphisms-lemma-quasi-compact-scheme-theoretic-image@cref}{{[originallemma][0][]01R8}{[1][0][]source}{}{}{}}
\makeatother
\crefname{originallemma}{original lemma}{original lemmas}
\Crefname{originallemma}{Original lemma}{Original lemmas}
\makeatletter
\newlabel{situation-setup}{{02QL}{source}{}{sourcepointer-8}{}}
\newlabel{situation-setup@cref}{{[originallemma][0][]02QL}{[1][0][]source}{}{}{}}
\makeatother
\title{Proper pushforward commutes with the k-cycle class of a coherent sheaf whose support has dimension at most k}
\author{The Stacks Project Authors}
\date{stacks-project:chow.tex:local-snapshot-20261012; GNU FDL 1.2 or later}
\begin{document}
\maketitle
\noindent\textit{Editorial note.} All delivered mathematical text is the author's own, copied from the pinned Stacks Project chapter file chow.tex; no statement or proof step was rewritten, reordered or removed.\par\medskip
\noindent\textit{Editorial note.} Preamble provenance: the pinned chapter file opens with \textbackslash{}input\{preamble\} and carries no preamble text of its own, so the project's single Stacks-wide preamble staged at sources/stacks-chow/preamble.tex is excerpted verbatim as two spans, lines 14-18 (Xy-pic setup) and lines 154-217 (the native theorem declarations and mathematical macros). Nothing was added, removed or redefined; the wrapper class and the amsmath/amssymb/amsthm/graphicx base packages are supplied by the assembler.\par\medskip
\noindent\textit{Editorial note.} Same-chapter pointers are located at their \textbackslash{}label lines inside the pinned chapter file; cross-chapter pointers are located in the project's own tag list staged in the snapshot (sources/stacks-chow/tags.txt), at the verbatim lines carrying their canonical tags.\par\medskip
\noindent\textit{Editorial note.} The author's parenthetical at source line 2040, 'The case of a closed immersion is straightforward (omitted).', is retained verbatim inside the counted proof; the closed-immersion case is the routine residue of an argument whose main reduction and final multiplicity comparison are fully carried out in the span.\par\medskip
\noindent\textit{Allowed standard prerequisites (editorial).}\begin{itemize}
\item Cycle class of a closed subscheme versus that of its structure sheaf. Same chapter chow.tex, label lemma-cycle-closed-coherent, invoked at proof line 2015. Reduces part (1) of the statement to part (2) applied to the structure sheaf O\_Z; part (1) is therefore not counted separately and is settled by the cited lemma plus the argument in the span.
\item Every coherent sheaf is the pushforward of a coherent sheaf on the closed subscheme given by its support. Cohomology of Schemes, labels coherent-lemma-coherent-support-closed (proof line 2018) and coherent-proposition-proper-pushforward-coherent (used in the statement at proof lines 2008-2011). Produces the closed immersion i : Z -> X and the coherent O\_Z-module G with i\_* G isomorphic to F, which is the first reduction of the proof.
\item Finiteness of the generic fibre and the pushdown formula for a finite ring map. Morphisms, labels morphisms-lemma-quasi-compact-scheme-theoretic-image (proof line 2042) and morphisms-lemma-generically-finite (proof line 2069); same chapter chow.tex, label lemma-multiplicity-finite (proof line 2066); Algebra, label algebra-lemma-pushdown-module (proof line 2094). Makes f|\_Z dominant and, after shrinking to an open neighbourhood of the generic point, finite and affine, and finally converts the coefficient of Z in f\_*[F]\_k into the sum over the primes of A over p via the pushdown formula.
\end{itemize}\medskip
\ifdefined\hypertarget\hypertarget{sourcepointer-1}{}\fi
\noindent\textit{Source reference.} Original lemma 02M0: sources/stacks-chow/tags.txt, source lines 3427--3427. Navigation pointer for \textbackslash{}ref\{algebra-lemma-pushdown-module\} used inside the counted span (line 2094); the argument that consumes it is retained in the span.\par\medskip
\ifdefined\hypertarget\hypertarget{sourcepointer-2}{}\fi
\noindent\textit{Source reference.} Original lemma 01Y5: sources/stacks-chow/tags.txt, source lines 2544--2544. Navigation pointer for \textbackslash{}ref\{coherent-lemma-coherent-support-closed\} used inside the counted span (line 2018); the argument that consumes it is retained in the span.\par\medskip
\ifdefined\hypertarget\hypertarget{sourcepointer-3}{}\fi
\noindent\textit{Source reference.} Original lemma 02O5: sources/stacks-chow/tags.txt, source lines 3504--3504. Navigation pointer for \textbackslash{}ref\{coherent-proposition-proper-pushforward-coherent\} used inside the counted span (line 2011); the argument that consumes it is retained in the span.\par\medskip
\ifdefined\hypertarget\hypertarget{sourcepointer-4}{}\fi
\noindent\textit{Source reference.} Original lemma 02QY: sources/stacks-chow/pinned-original-source.tex, source lines 1776--1776. Navigation pointer for \textbackslash{}ref\{lemma-cycle-closed-coherent\} used inside the counted span (line 2015); the argument that consumes it is retained in the span.\par\medskip
\ifdefined\hypertarget\hypertarget{sourcepointer-5}{}\fi
\noindent\textit{Source reference.} Original lemma 02QT: sources/stacks-chow/pinned-original-source.tex, source lines 1629--1629. Navigation pointer for \textbackslash{}ref\{lemma-multiplicity-finite\} used inside the counted span (line 2066); the argument that consumes it is retained in the span.\par\medskip
\ifdefined\hypertarget\hypertarget{sourcepointer-6}{}\fi
\noindent\textit{Source reference.} Original lemma 02NW: sources/stacks-chow/tags.txt, source lines 3495--3495. Navigation pointer for \textbackslash{}ref\{morphisms-lemma-generically-finite\} used inside the counted span (line 2069); the argument that consumes it is retained in the span.\par\medskip
\ifdefined\hypertarget\hypertarget{sourcepointer-7}{}\fi
\noindent\textit{Source reference.} Original lemma 01R8: sources/stacks-chow/tags.txt, source lines 2293--2293. Navigation pointer for \textbackslash{}ref\{morphisms-lemma-quasi-compact-scheme-theoretic-image\} used inside the counted span (line 2042); the argument that consumes it is retained in the span.\par\medskip
\ifdefined\hypertarget\hypertarget{sourcepointer-8}{}\fi
\noindent\textit{Source reference.} Original lemma 02QL: sources/stacks-chow/pinned-original-source.tex, source lines 1372--1372. Navigation pointer for \textbackslash{}ref\{situation-setup\} used inside the counted span (line 1993); the argument that consumes it is retained in the span.\par\medskip
\begin{lemma}
\label{lemma-cycle-push-sheaf}
Let $(S, \delta)$ be as in Situation \ref{situation-setup}.
Let $f : X \to Y$ be a proper morphism of schemes which are
locally of finite type over $S$.
\begin{enumerate}
\item Let $Z \subset X$ be a closed subscheme with $\dim_\delta(Z) \leq k$.
Then
$$
f_*[Z]_k = [f_*{\mathcal O}_Z]_k.
$$
\item Let $\mathcal{F}$ be a coherent sheaf on $X$ such that
$\dim_\delta(\text{Supp}(\mathcal{F})) \leq k$. Then
$$
f_*[\mathcal{F}]_k = [f_*{\mathcal F}]_k.
$$
\end{enumerate}
Note that the statement makes sense since $f_*\mathcal{F}$ and
$f_*\mathcal{O}_Z$ are coherent $\mathcal{O}_Y$-modules by
Cohomology of Schemes, Proposition
\ref{coherent-proposition-proper-pushforward-coherent}.
\end{lemma}

\begin{proof}
Part (1) follows from (2) and Lemma \ref{lemma-cycle-closed-coherent}.
Let $\mathcal{F}$ be a coherent sheaf on $X$.
Assume that $\dim_\delta(\text{Supp}(\mathcal{F})) \leq k$.
By Cohomology of Schemes, Lemma \ref{coherent-lemma-coherent-support-closed}
there exists a closed subscheme $i : Z \to X$ and a coherent
$\mathcal{O}_Z$-module $\mathcal{G}$ such that
$i_*\mathcal{G} \cong \mathcal{F}$ and such that the support
of $\mathcal{F}$ is $Z$. Let $Z' \subset Y$ be the scheme theoretic image
of $f|_Z : Z \to Y$. Consider the commutative diagram of schemes
$$
\xymatrix{
Z \ar[r]_i \ar[d]_{f|_Z} &
X \ar[d]^f \\
Z' \ar[r]^{i'} & Y
}
$$
We have $f_*\mathcal{F} = f_*i_*\mathcal{G} = i'_*(f|_Z)_*\mathcal{G}$
by going around the diagram in two ways. Suppose we know the result holds
for closed immersions and for $f|_Z$. Then we see that
$$
f_*[\mathcal{F}]_k = f_*i_*[\mathcal{G}]_k
= (i')_*(f|_Z)_*[\mathcal{G}]_k =
(i')_*[(f|_Z)_*\mathcal{G}]_k =
[(i')_*(f|_Z)_*\mathcal{G}]_k = [f_*\mathcal{F}]_k
$$
as desired. The case of a closed immersion is straightforward (omitted).
Note that $f|_Z : Z \to Z'$ is a dominant morphism (see
Morphisms, Lemma \ref{morphisms-lemma-quasi-compact-scheme-theoretic-image}).
Thus we have reduced to the case where
$\dim_\delta(X) \leq k$ and $f : X \to Y$ is proper and dominant.

\medskip\noindent
Assume $\dim_\delta(X) \leq k$ and $f : X \to Y$ is proper and dominant.
Since $f$ is dominant, for every irreducible component $Z \subset Y$
with generic point $\eta$ there exists a point $\xi \in X$ such
that $f(\xi) = \eta$. Hence $\delta(\eta) \leq \delta(\xi) \leq k$.
Thus we see that in the expressions
$$
f_*[\mathcal{F}]_k = \sum n_Z[Z],
\quad
\text{and}
\quad
[f_*\mathcal{F}]_k = \sum m_Z[Z].
$$
whenever $n_Z \not = 0$, or $m_Z \not = 0$ the integral closed
subscheme $Z$ is actually an irreducible component of $Y$ of
$\delta$-dimension $k$. Pick such an integral closed subscheme
$Z \subset Y$ and denote $\eta$ its generic point. Note that for
any $\xi \in X$ with $f(\xi) = \eta$ we have $\delta(\xi) \geq k$
and hence $\xi$ is a generic point of an irreducible component
of $X$ of $\delta$-dimension $k$ as well
(see Lemma \ref{lemma-multiplicity-finite}). Since $f$ is quasi-compact
and $X$ is locally Noetherian, there can be only finitely many of
these and hence $f^{-1}(\{\eta\})$ is finite.
By Morphisms, Lemma \ref{morphisms-lemma-generically-finite} there exists
an open neighbourhood $\eta \in V \subset Y$ such that $f^{-1}(V) \to V$
is finite. Replacing $Y$ by $V$ and $X$ by $f^{-1}(V)$ we reduce to the
case where $Y$ is affine, and $f$ is finite.

\medskip\noindent
Write $Y = \Spec(R)$ and $X = \Spec(A)$ (possible as
a finite morphism is affine).
Then $R$ and $A$ are Noetherian rings and $A$ is finite over $R$.
Moreover $\mathcal{F} = \widetilde{M}$ for some finite $A$-module
$M$. Note that $f_*\mathcal{F}$ corresponds to $M$ viewed as an $R$-module.
Let $\mathfrak p \subset R$ be the minimal prime corresponding
to $\eta \in Y$. The coefficient of $Z$ in $[f_*\mathcal{F}]_k$
is clearly $\text{length}_{R_{\mathfrak p}}(M_{\mathfrak p})$.
Let $\mathfrak q_i$, $i = 1, \ldots, t$ be the primes of $A$
lying over $\mathfrak p$. Then $A_{\mathfrak p} = \prod A_{\mathfrak q_i}$
since $A_{\mathfrak p}$ is an Artinian ring being finite over the
dimension zero local Noetherian ring $R_{\mathfrak p}$.
Clearly the coefficient of $Z$ in $f_*[\mathcal{F}]_k$ is
$$
\sum\nolimits_{i = 1, \ldots, t}
[\kappa(\mathfrak q_i) : \kappa(\mathfrak p)]
\text{length}_{A_{\mathfrak q_i}}(M_{\mathfrak q_i})
$$
Hence the desired equality follows from
Algebra, Lemma \ref{algebra-lemma-pushdown-module}.
\end{proof}

\end{document}
